\documentclass[11pt,a4paper]{article}
\usepackage[T1]{fontenc}
\usepackage[utf8]{inputenc}
\usepackage[polish]{babel}
\usepackage{lmodern}
\usepackage[a4paper,margin=2.4cm]{geometry}
\usepackage{amsmath}
\usepackage{hyperref}

\title{HF20.2: jawna obsługa dostępności MLflow}
\author{Edward Szczypka}
\date{\today}

\begin{document}
\maketitle

\section{Problem}

Samo ustawienie \texttt{mlflow.enabled=true} nie dowodzi, że serwer pod
\texttt{tracking\_uri} jest uruchomiony. Dotychczas błąd połączenia był
ujawniany dopiero podczas inicjalizacji eksperymentu i powodował długi
traceback.

\section{Preflight}

Przed treningiem definiujemy stan:
\[
S\in\{
\text{ready},
\text{disabled},
\text{not-installed},
\text{not-running},
\text{invalid}
\}.
\]

Trening z wymaganym MLflow może rozpocząć się tylko dla
\(S=\text{ready}\). Dla treningu bez MLflow stosowane jest jednorazowe
nadpisanie środowiska:
\[
\texttt{SMOG\_AI\_MLFLOW\_ENABLED=false}.
\]

\section{Polityki}

\begin{description}
\item[required] MLflow musi być dostępny; w przeciwnym razie trening nie
rozpoczyna się.
\item[disabled] trening zawsze działa bez MLflow.
\item[auto] używa MLflow, gdy jest dostępny, w przeciwnym razie kontynuuje
bez niego z jawnym ostrzeżeniem.
\item[ask] użytkownik wybiera ponowny test, trening bez MLflow albo
zakończenie.
\end{description}

\end{document}
